% GENERATED FILE. Edit canonical lessons, then run npm run generate:books.
% canonical-source-hash: fnv1a64:3b50e0a7d312e307
% canonical-lessons: KA-C507-tegu, KA-C507-mitukisu, KA-C507-edusiru-bidu, KA-C507-bikkalike, KA-C507-ubbasa

\chapter{Small Body Sounds}
\label{ch:ka-small-body-sounds}
\hlchaptermodality{507}

\begin{chapteropening}
\textbf{By the end of this chapter:} \emph{I can say to belch, to blink, to pant, hiccups and asthma.}

Complete the last lesson of chapter 507: I can say to belch, to blink, to pant, hiccups and asthma.
\end{chapteropening}

\section[\texorpdfstring{\kn{ತೇಗು} (tēgu)}{ತೇಗು (tēgu)}]{\kn{ತೇಗು} (tēgu) — to belch; to burp}

\label{lesson:KA-C507-tegu}



\begin{quote}
Before the new one: say the Kannada for to stretch, then the Kannada for an eyelid.
\end{quote}

\subsection*{You'll want to know: \kn{ತೇಗು}}
\textbf{\kn{ತೇಗು}} — \emph{tēgu} — \textquotedblleft{}to belch; to burp\textquotedblright{}.

\begin{quote}
\kn{ಊಟದ} \kn{ನಂತರ} \kn{ಮಗು} \kn{ತೇಗಿತು}.
\end{quote}

\begin{quote}
\emph{ūṭada nantara magu tēgitu.}
\end{quote}

\begin{quote}
After the meal the baby burped.
\end{quote}

Past (I …): \kn{ತೇಗಿದೆ}, \emph{tēgide}.

\kn{ತೇಗು} is also the burp itself.

\begin{cousinweb}
\textbf{\kn{ತೇಗು}} (\emph{tēgu}) is native Dravidian: DEDR 3451 lists Kannada \emph{tēgu}, \emph{tēku}, \textquotedblleft{}to belch; a belch\textquotedblright{}, beside Tamil \emph{tēkku}.
\end{cousinweb}

\begin{grammarlens}[title={What you've built}]
One more word for describing what you lived and felt.
\end{grammarlens}

\subsection*{Your turn}
\begin{itemize}
\raggedright
  \item \emph{Say it:} \emph{tēgu}
  \item \emph{Say it:} \emph{tēgu}, once more
  \item \emph{Say it:} \emph{ūṭada nantara magu tēgitu.} — After the meal the baby burped.
  \item \emph{Recall:} say the Kannada for to stretch, then the Kannada for an eyelid, then say \emph{tēgu} again
\end{itemize}

\begin{tcolorbox}[breakable,colback=teal!4,colframe=teal!35!black,title={Before you move on}]
What does \kn{ತೇಗು} mean? (\textquotedblleft{}To belch; to burp\textquotedblright{}.) Say it once more.
\end{tcolorbox}
\section[\texorpdfstring{\kn{ಮಿಟುಕಿಸು} (miṭukisu)}{ಮಿಟುಕಿಸು (miṭukisu)}]{\kn{ಮಿಟುಕಿಸು} (miṭukisu) — to blink; to wink}

\label{lesson:KA-C507-mitukisu}



\begin{quote}
Before the new one: say the Kannada for an eyelid, then the Kannada for to belch.
\end{quote}

\subsection*{You'll want to know: \kn{ಮಿಟುಕಿಸು}}
\textbf{\kn{ಮಿಟುಕಿಸು}} — \emph{miṭukisu} — \textquotedblleft{}to blink; to wink\textquotedblright{}.

\begin{quote}
\kn{ಬಿಸಿಲಿನಲ್ಲಿ} \kn{ಅವನು} \kn{ಕಣ್ಣು} \kn{ಮಿಟುಕಿಸಿದನು}.
\end{quote}

\begin{quote}
\emph{bisilinalli avanu kaṇṇu miṭukisidanu.}
\end{quote}

\begin{quote}
In the sunshine he blinked.
\end{quote}

Past (I …): \kn{ಮಿಟುಕಿಸಿದೆ}, \emph{miṭukiside}.

\kn{ಕಣ್ಣು} \kn{ಮಿಟುಕಿಸು}, to blink; to wink at someone, add -\kn{ಗೆ}.

\begin{cousinweb}
\textbf{\kn{ಮಿಟುಕಿಸು}} (\emph{miṭukisu}) is native Dravidian: DEDR 4845 lists Kannada \emph{miṭakisu}, \textquotedblleft{}to open and shut the eyes rapidly, blink, wink\textquotedblright{}.
\end{cousinweb}

\begin{grammarlens}[title={What you've built}]
One more word for describing what you lived and felt.
\end{grammarlens}

\subsection*{Your turn}
\begin{itemize}
\raggedright
  \item \emph{Say it:} \emph{miṭukisu}
  \item \emph{Say it:} \emph{miṭukisu}, once more
  \item \emph{Say it:} \emph{bisilinalli avanu kaṇṇu miṭukisidanu.} — In the sunshine he blinked.
  \item \emph{Recall:} say the Kannada for an eyelid, then the Kannada for to belch, then say \emph{miṭukisu} again
\end{itemize}

\begin{tcolorbox}[breakable,colback=teal!4,colframe=teal!35!black,title={Before you move on}]
What does \kn{ಮಿಟುಕಿಸು} mean? (\textquotedblleft{}To blink; to wink\textquotedblright{}.) Say it once more.
\end{tcolorbox}
\section[\texorpdfstring{\kn{ಏದುಸಿರು} \kn{ಬಿಡು} (ēdusiru biḍu)}{ಏದುಸಿರು ಬಿಡು (ēdusiru biḍu)}]{\kn{ಏದುಸಿರು} \kn{ಬಿಡು} (ēdusiru biḍu) — to pant; to gasp}

\label{lesson:KA-C507-edusiru-bidu}



\begin{quote}
Before the new one: say the Kannada for to belch, then the Kannada for to blink.
\end{quote}

\subsection*{You'll want to know: \kn{ಏದುಸಿರು} \kn{ಬಿಡು}}
\textbf{\kn{ಏದುಸಿರು} \kn{ಬಿಡು}} — \emph{ēdusiru biḍu} — \textquotedblleft{}to pant; to gasp\textquotedblright{}.

\begin{quote}
\kn{ಮೆಟ್ಟಿಲು} \kn{ಹತ್ತಿ} \kn{ಅಜ್ಜಿ} \kn{ಏದುಸಿರು} \kn{ಬಿಟ್ಟರು}.
\end{quote}

\begin{quote}
\emph{meṭṭilu hatti ajji ēdusiru biṭṭaru.}
\end{quote}

\begin{quote}
Grandmother panted after climbing the stairs.
\end{quote}

Past (I …): \kn{ಏದುಸಿರು} \kn{ಬಿಟ್ಟೆ}, \emph{ēdusiru biṭṭe}.

\kn{ಏದುಸಿರು} is panting breath, built on \kn{ಉಸಿರು}, breath; \kn{ಬಿಡು}, to let out.

\begin{cousinweb}
\textbf{\kn{ಏದುಸಿರು} \kn{ಬಿಡು}} (\emph{ēdusiru biḍu}) is a native Kannada phrase: \emph{usiru}, \textquotedblleft{}breath\textquotedblright{} (DEDR 645), and \emph{biḍu}, \textquotedblleft{}to let go\textquotedblright{} (DEDR 5393). The origin of the first part, ēd-, is not settled.
\end{cousinweb}

\begin{grammarlens}[title={What you've built}]
One more word for describing what you lived and felt.
\end{grammarlens}

\subsection*{Your turn}
\begin{itemize}
\raggedright
  \item \emph{Say it:} \emph{ēdusiru biḍu}
  \item \emph{Say it:} \emph{ēdusiru biḍu}, once more
  \item \emph{Say it:} \emph{meṭṭilu hatti ajji ēdusiru biṭṭaru.} — Grandmother panted after climbing the stairs.
  \item \emph{Recall:} say the Kannada for to belch, then the Kannada for to blink, then say \emph{ēdusiru biḍu} again
\end{itemize}

\begin{tcolorbox}[breakable,colback=teal!4,colframe=teal!35!black,title={Before you move on}]
What does \kn{ಏದುಸಿರು} \kn{ಬಿಡು} mean? (\textquotedblleft{}To pant; to gasp\textquotedblright{}.) Say it once more.
\end{tcolorbox}
\section[\texorpdfstring{\kn{ಬಿಕ್ಕಳಿಕೆ} (bikkaḷike)}{ಬಿಕ್ಕಳಿಕೆ (bikkaḷike)}]{\kn{ಬಿಕ್ಕಳಿಕೆ} (bikkaḷike) — hiccups}

\label{lesson:KA-C507-bikkalike}



\begin{quote}
Before the new one: say the Kannada for to blink, then the Kannada for to pant.
\end{quote}

\subsection*{You'll want to know: \kn{ಬಿಕ್ಕಳಿಕೆ}}
\textbf{\kn{ಬಿಕ್ಕಳಿಕೆ}} — \emph{bikkaḷike} — \textquotedblleft{}hiccups\textquotedblright{}.

\begin{quote}
\kn{ಬೇಗ} \kn{ಬೇಗ} \kn{ತಿಂದಿದ್ದರಿಂದ} \kn{ನನಗೆ} \kn{ಬಿಕ್ಕಳಿಕೆ} \kn{ಬಂದಿತು}.
\end{quote}

\begin{quote}
\emph{bēga bēga tindiddarinda nanage bikkaḷike banditu.}
\end{quote}

\begin{quote}
Because I ate too fast, I got hiccups.
\end{quote}

Hiccups \textquotedblleft{}come\textquotedblright{}: \kn{ಬಿಕ್ಕಳಿಕೆ} \kn{ಬರು}.

\begin{cousinweb}
\textbf{\kn{ಬಿಕ್ಕಳಿಕೆ}} (\emph{bikkaḷike}) is native Dravidian: DEDR 5383 sets Kannada \emph{bikkaḷike} beside Tamil \emph{vikku}, \textquotedblleft{}to hiccup\textquotedblright{}.
\end{cousinweb}

\begin{grammarlens}[title={What you've built}]
One more word for describing what you lived and felt.
\end{grammarlens}

\subsection*{Your turn}
\begin{itemize}
\raggedright
  \item \emph{Say it:} \emph{bikkaḷike}
  \item \emph{Say it:} \emph{bikkaḷike}, once more
  \item \emph{Say it:} \emph{bēga bēga tindiddarinda nanage bikkaḷike banditu.} — Because I ate too fast, I got hiccups.
  \item \emph{Recall:} say the Kannada for to blink, then the Kannada for to pant, then say \emph{bikkaḷike} again
\end{itemize}

\begin{tcolorbox}[breakable,colback=teal!4,colframe=teal!35!black,title={Before you move on}]
What does \kn{ಬಿಕ್ಕಳಿಕೆ} mean? (\textquotedblleft{}Hiccups\textquotedblright{}.) Say it once more.
\end{tcolorbox}
\section[\texorpdfstring{\kn{ಉಬ್ಬಸ} (ubbasa)}{ಉಬ್ಬಸ (ubbasa)}]{\kn{ಉಬ್ಬಸ} (ubbasa) — asthma; breathlessness}

\label{lesson:KA-C507-ubbasa}



\begin{quote}
Before the new one: say the Kannada for to pant, then the Kannada for hiccups.
\end{quote}

\subsection*{You'll want to know: \kn{ಉಬ್ಬಸ}}
\textbf{\kn{ಉಬ್ಬಸ}} — \emph{ubbasa} — \textquotedblleft{}asthma; breathlessness\textquotedblright{}.

\begin{quote}
\kn{ಚಳಿಗಾಲದಲ್ಲಿ} \kn{ಅಜ್ಜನ} \kn{ಉಬ್ಬಸ} \kn{ಹೆಚ್ಚಾಗುತ್ತದೆ}.
\end{quote}

\begin{quote}
\emph{caḷigāladalli ajjana ubbasa heccāguttade.}
\end{quote}

\begin{quote}
In winter grandfather's asthma gets worse.
\end{quote}

Doctors also say \kn{ಆಸ್ತಮಾ}.

\begin{cousinweb}
\textbf{\kn{ಉಬ್ಬಸ}} (\emph{ubbasa}) is native Dravidian: DEDR 632 lists Kannada \emph{ubbasa} with Tamil and Kota words for asthma and panting.
\end{cousinweb}

\begin{grammarlens}[title={What you've built}]
One more word for describing what you lived and felt.
\end{grammarlens}

\subsection*{Your turn}
\begin{itemize}
\raggedright
  \item \emph{Say it:} \emph{ubbasa}
  \item \emph{Say it:} \emph{ubbasa}, once more
  \item \emph{Say it:} \emph{caḷigāladalli ajjana ubbasa heccāguttade.} — In winter grandfather's asthma gets worse.
  \item \emph{Recall:} say the Kannada for to pant, then the Kannada for hiccups, then say \emph{ubbasa} again
\end{itemize}

\begin{tcolorbox}[breakable,colback=teal!4,colframe=teal!35!black,title={Before you move on}]
What does \kn{ಉಬ್ಬಸ} mean? (\textquotedblleft{}Asthma; breathlessness\textquotedblright{}.) Say it once more.
\end{tcolorbox}
